\section{归纳集}

\begin{definition}{归纳集}{}
	设$E$是偏序集.如果$E$的每个是全序子集在$E$中都有上确界,则称$E$是归纳集.
\end{definition}

\begin{example}{}{}
	\begin{enumerate}
		\item 设$A$是集合,$\mathfrak{F}\subseteq\mathcal{P}\left(A\right)$,配备包含偏序.如果$\mathfrak{F}$的每个全序子集$\mathfrak{O}$满足$\bigcup\limits_{O\in\mathfrak{O}}O\in\mathfrak{F}$,则$\mathfrak{F}$是归纳集.
		\item 设$A,B$是集合,给$\PFunc\left(A,B\right)$配备函数延拓偏序后它是归纳集.
		\item 根据无穷公理,$\N$在标准全序下不是归纳集,因为$\N$在这个全序下没有上界.
	\end{enumerate}
\end{example}

\begin{proposition}{加强版的Zorn引理}{}
	设$E$是偏序集.如果$E$的每个良序子集在$E$中都有上界,则$E$有极大元.
\end{proposition}

\begin{theorem}{Zorn引理}{}
	每个归纳集都有极大元.
\end{theorem}

\begin{corollary}{}{}
	设归纳集$E$配备偏序$\leq$,$a\in E$,则$E$中有一个极大元$m$满足$m\geq a$.
\end{corollary}

\begin{corollary}{}{}
	设$E$是集合,$\mathfrak{F}\subseteq\mathcal{P}\left(E\right)$.
	\begin{enumerate}
		\item 若对$\mathfrak{F}$的每个全序子集$\mathfrak{P}$都有$\bigcup\limits_{P\in\mathfrak{P}}P\in\mathfrak{F}$,则$\mathfrak{F}$有极大元.
		\item 若对$\mathfrak{F}$的每个全序子集$\mathfrak{P}$都有$\bigcap\limits_{P\in\mathfrak{P}}P\in\mathfrak{F}$,则$\mathfrak{F}$有极小元.
	\end{enumerate}
\end{corollary}













































